\documentclass[12pt]{article}

\parskip 1ex plus 0.4ex minus 0.4ex
\parindent0ex

\usepackage{enumerate}
\usepackage{prettyref}
\usepackage{pdfpages}
\usepackage{float}
\usepackage{dirtree}
\usepackage{xltxtra}
\usepackage{polyglossia}
\setdefaultlanguage[spelling=new]{german}
\usepackage{fontspec}
\usepackage{listings}
\usepackage{xcolor}
\usepackage{amsmath}
\usepackage{amsthm}
\usepackage{amstext}
\usepackage{amssymb}
\usepackage{graphicx}
\usepackage[colorlinks,
pdfpagelabels,
pdfstartview = FitH,
bookmarksopen = true,
bookmarksnumbered = true,
linkcolor = blue, %Farbe
plainpages = false,
hypertexnames = false,
citecolor = black,
xetex] {hyperref}
\usepackage[style=alphabetic,backend=biber]{biblatex}
\bibliography{bib}

\title{Praktikumsbericht: Einf\"uhrung in Python}
\author{Marcus Ganske, 36603\\
		Lukas Krieg, 53506}
\date{\today}

\definecolor{light-gray}{gray}{0.95}

\lstset{
	language=python,
	breaklines,
	basicstyle=\ttfamily\small,
	backgroundcolor=\color{light-gray},
	keywordstyle=\color{blue},
	stringstyle=\color{olive},
	commentstyle=\color{gray}\ttfamily,
	numbers = left,
	numberstyle = \tiny,
	numberblanklines=true,
	stepnumber = 1,
	tabsize = 4,
	numbersep=10pt,
	xleftmargin=10pt,
}

\begin{document}
\maketitle
\vspace{+8cm}{
}
\includegraphics[width=12cm]{Hochschule-aalen.pdf}

\newpage
%Inhaltsverzeichnis
\renewcommand\contentsname{Inhaltsverzeichnis}
\tableofcontents
\newpage
	
	\section{Einleitung}
	\begin{figure}[H]
			\dirtree{%Baumstruktur
				.1 Code/.
				.2 Aufgabe1.c.
				.2 shellcode.asm.
				.2 shellcode.bin.
				.2 shellcode.o.
				.2 shellcode.py.
			}
	
		\caption{Aufbau des Ordners Code mit allen Dateien}
	\end{figure}	


\newpage
\section{Aufgabe 1: Codeanalyse \& Reverse Engineering}
In diesem Kapitel wird auf Aufgabe 1 des Praktikums eingegangen.
Es sollte der folgende Code implementiert und analysiert werden.
\begin{lstlisting}
include <unistd.h>
int main() {
	char *cmd[] =  {"/bin/sh","-c","ls-l", (char*)0 };
	int ret;
	ret = execve(cmd[0], cmd, NULL);
	return ret;
}
\end{lstlisting}

Der Quellcode des C Programms wurde in der Datei unistd.c implementiert.
Diese Datei befindet sich im verzeichnis Code.

\subsection{Codeanalyse}
\textbf{Beschreiben sie die Funktionsweise des Codes}\\
Das C-Programm initialisiert in Zeile 3 ein String Array mit Namen cmd der Inhalt der Strings ist: \\
\begin{lstlisting}
"/bin/sh"
"-c"
"ls -l"
NULL Pointer Value von Typ char.
\end{lstlisting}
Danach wird in Zeile 4 ein integer mit Variablennamen "ret" initialisiert.\\
Die Funktion execve wird aufgerufen mit den Übergabeparametern.
\begin{lstlisting}
cmd[0] 	-> "/bin/sh"
cmd		-> "/bin/sh", "-c", "ls -l", NULL Pointer Value
NULL	 
\end{lstlisting}

\textbf{Funktionsaufruf Execve}

Die Funktion Execve führt Programme aus, auf welche mit einem Dateipfad gezeigt wird.
In diesem Fall handelt es sich um das erste Übergabeparameter cmd[0]
worin sich der Pfad "/bin/bash" befindet. Es wird eine Shell gestartet.


Das zweite Übergabeparameter von Execve muss ein Array aus Argumentstrings sein
die an das neue Programm übergeben werden. In diesem Fall:

\begin{lstlisting}
cmd = {"/bin/sh","-c","ls -l",(char*)0}
\end{lstlisting}
das erste Argument "/bin/sh" gibt den Interpreter  an
das zweite Argument "-c" bedeutet, dass das folgende command ausgeführt werden soll
das dritte Argument "ls -l" wird durch die option "-c" ausgeführt
das vierte Argument (char*) 0 ist ein Nullpointer, der benötigt wird
um das Ende des Arrays für execve ersichtlich zu machen

Das dritte Übergabeparameter von execve ist "NULL".
Hier wird ein Stringarray mit zusätzlichen Variablen erwartet, die an
das neue Programm weitergereicht werden können.
Dieses Parameter ist genau wie das Zweite nullterminiert.
Deswegen wird hier einfach NULL übergeben.n


Der Returnwert wird in die integervariable "ret" geschrieben.
Bei Erfolg gibt es bei  execve allerdings keinen Rückgabewert.
Bei einem Fehler wird allerdings -1 zurückgegeben.

\subsection{Reverse Engineering mit gdb}
\begin{lstlisting}
Wie werden die Übergabeparameter an die Funktion
übergeben 
Die übergabeparameter werden in die Register 
rdx,rdi,rsi gespeichert

Bild Register rdx rdi rsi vor call

Im Register rdi befindet sich durch 
Zeile 
main+80   mov rdx, rax
die Adresse 0x555555554814
des Strings "/bin/sh"
 

Zeile main+77 mov rsi, rcx 
Im Register rsi befindet sich die Adresse
des ersten Strings unserers Stringarrays cmd.
Die Adresse lautet 0x7fffffffde20. 
Diese Adresse liegt im Stack. 
Der Stack sieht vor dem Aufruf der Funktion folgendermaßen aus.
0x7fffffffde10:	0x1	0x5555555547dd <__libc_csu_init+77>
0x7fffffffde20:	0x555555554814	0x55555555481c
0x7fffffffde30:	0x55555555481f	0x0
0x7fffffffde40:	0x7fffffffdf30	0xc7f5fc08a85bb400
0x7fffffffde50:	0x555555554790 <__libc_csu_init>	0x7ffff7a313f1 <__libc_start_main+241>

Die auf  0x7fffffffde20 folgenden Adressen enthalten die benötigten Strings
0x7fffffffde20: 0x555555554814  -> "/bin/sh"
0x7fffffffde28: 0x55555555481c	-> "-c"
0x7fffffffde30: 0x55555555481f	-> "ls -l"
0x7fffffffde38: 0x0		-> 0
somit wird über den register rsi die Variable cmd übergeben.
mov eax, 0x0 
ein NULL, welches das dritte Übergabeparameter 
von execve ist.


lea rax,[rip+0xe6] 
Die Adressen der Strings aus dem 
Teil des Rams in der der Programmcode 
steht in den register rax geschrieben 

mov QWORD PTR [rbp-0x30],rax
danach wird die Adresse aus dem rax in 
den Stack geschrieben

ii) Wie und an welcher Stelle werden die übergebenen Parameter im Speicher abgelegt?
Dump of assembler code for function main:
   0x0000555555554710 <+0>:	push   rbp
   0x0000555555554711 <+1>:	mov    rbp,rsp
   0x0000555555554714 <+4>:	sub    rsp,0x40
   0x0000555555554718 <+8>:	mov    rax,QWORD PTR fs:0x28
   0x0000555555554721 <+17>:	mov    QWORD PTR [rbp-0x8],rax
   0x0000555555554725 <+21>:	xor    eax,eax
   0x0000555555554727 <+23>:	lea    rax,[rip+0xe6]        # 0x555555554814
   0x000055555555472e <+30>:	mov    QWORD PTR [rbp-0x30],rax
   0x0000555555554732 <+34>:	lea    rax,[rip+0xe3]        # 0x55555555481c
   0x0000555555554739 <+41>:	mov    QWORD PTR [rbp-0x28],rax
   0x000055555555473d <+45>:	lea    rax,[rip+0xdb]        # 0x55555555481f
   0x0000555555554744 <+52>:	mov    QWORD PTR [rbp-0x20],rax
   0x0000555555554748 <+56>:	mov    QWORD PTR [rbp-0x18],0x0
   0x0000555555554750 <+64>:	mov    rax,QWORD PTR [rbp-0x30]
   0x0000555555554754 <+68>:	lea    rcx,[rbp-0x30]
   0x0000555555554758 <+72>:	mov    edx,0x0
   0x000055555555475d <+77>:	mov    rsi,rcx
   0x0000555555554760 <+80>:	mov    rdi,rax
   0x0000555555554763 <+83>:	call   0x5555555545c8
   0x0000555555554768 <+88>:	mov    DWORD PTR [rbp-0x34],eax
   0x000055555555476b <+91>:	mov    eax,DWORD PTR [rbp-0x34]
   0x000055555555476e <+94>:	mov    rdx,QWORD PTR [rbp-0x8]
   0x0000555555554772 <+98>:	xor    rdx,QWORD PTR fs:0x28
   0x000055555555477b <+107>:	je     0x555555554782 <main+114>
   0x000055555555477d <+109>:	call   0x5555555545c0
   0x0000555555554782 <+114>:	leave  
   0x0000555555554783 <+115>:	ret  

+23 Adresse von "/bin/sh" wird aus Programmcode in Register rax geschrieben 
+30 Adresse von "/bin/sh" wird aus rax auf den stack geschrieben
+34 Adresse von "-c" wird aus Programmcode in rax geschrieben
+41 Adresse von "-c" wird aus Programmcode in rax geschrieben
+45 Adresse von "ls -l" wird aus Programmcode in rax geschrieben
+52 Adresse von "ls -l" wird aus Register rax auf den Stack geschrieben
+56 Adresse xx auf dem Stack wird mit 0 beschrieben.
Im Stack werden die Stringadressen als QWORD Pointer gespeichert, d.h 8byte pointer =64bit

Danach werden die Parameter wie oben beschrieben aus dem Stack wieder in Register geschrieben 
um nach dem Funktionsaufruf verwendet werden zu können.


iii)
Im stack befindet sich die Rücksprungadresse auf die mainfunktion
(gdb) x/a $rsp
0x7fffffffde08:	0x555555554768 <main+88>

0x555555554768
ist die adresse die der rip nach ausführen der funktion execve wieder aufnehmen muss

dies ist der wiedereintrittspunkt in der mainfunktion

0x0000555555554763 <+83>:	call   0x5555555545c8
0x0000555555554768 <+88>:	mov    DWORD PTR [rbp-0x34],eax
0x000055555555476b <+91>:	mov    eax,DWORD PTR [rbp-0x34]
0x000055555555476e <+94>:	mov    rdx,QWORD PTR [rbp-0x8]
0x0000555555554772 <+98>:	xor    rdx,QWORD PTR fs:0x28
0x000055555555477b <+107>:	je     0x555555554782 <main+114>
0x000055555555477d <+109>:	call   0x5555555545c0
0x0000555555554782 <+114>:	leave  
0x0000555555554783 <+115>:	ret    
\end{lstlisting}
\subsection{to-do}
BILDER NOCH EINFÜGEN!!!
FORMATIEREN





\end{document}


